\documentclass[11pt,a4paper]{article}

% --- Codificación e idioma ---
\usepackage[utf8]{inputenc}
\usepackage[T1]{fontenc}
\usepackage[spanish,es-noquoting]{babel}
\usepackage{lmodern}

% --- Página y tipografía ---
\usepackage[a4paper,margin=2.4cm]{geometry}
\usepackage{microtype}
\usepackage[htt]{hyphenat}
\usepackage{enumitem}
\setlist{leftmargin=1.4em}
\linespread{1.04}
\setlength{\emergencystretch}{3em}

% --- Matemática ---
\usepackage{amsmath,amssymb}

% --- Color y código ---
\usepackage{xcolor}
\definecolor{azulUNI}{RGB}{22,50,79}
\definecolor{grisfondo}{RGB}{246,248,250}
\definecolor{comentario}{RGB}{90,110,130}
\definecolor{cadena}{RGB}{140,70,40}
\usepackage{listings}
\lstdefinestyle{py}{
  language=Python,
  basicstyle=\ttfamily\footnotesize,
  backgroundcolor=\color{grisfondo},
  commentstyle=\color{comentario}\itshape,
  keywordstyle=\color{azulUNI}\bfseries,
  stringstyle=\color{cadena},
  frame=single, rulecolor=\color{gray!40},
  breaklines=true, columns=flexible, showstringspaces=false,
  xleftmargin=6pt, xrightmargin=6pt, aboveskip=5pt, belowskip=5pt,
  literate={á}{{\'a}}1 {é}{{\'e}}1 {í}{{\'i}}1 {ó}{{\'o}}1 {ú}{{\'u}}1 {ñ}{{\~n}}1 {Á}{{\'A}}1
}
\lstset{style=py}

% --- Diagramas ---
\usepackage{tikz}
\usetikzlibrary{arrows.meta,calc,positioning,fit,backgrounds}

% --- Encabezados ---
\usepackage{titlesec}
\titleformat{\section}{\normalfont\Large\bfseries\color{azulUNI}}{\thesection}{0.6em}{}
\titleformat{\subsection}{\normalfont\large\bfseries\color{azulUNI!85!black}}{\thesubsection}{0.6em}{}
\usepackage{fancyhdr}
\pagestyle{fancy}\fancyhf{}
\fancyhead[L]{\footnotesize Criptografía · Actividad III}
\fancyhead[R]{\footnotesize Custodia de claves (Modelo A)}
\fancyfoot[C]{\thepage}
\renewcommand{\headrulewidth}{0.4pt}

\usepackage[hidelinks]{hyperref}

% Estilos de diagrama de secuencia
\tikzset{
  actor/.style={draw=azulUNI,fill=azulUNI!12,rounded corners=2pt,
    minimum width=2.9cm,minimum height=0.75cm,font=\small\bfseries,align=center},
  msg/.style={-{Latex[length=2mm]},thick},
  rmsg/.style={-{Latex[length=2mm]},thick,densely dashed},
  lbl/.style={font=\scriptsize,midway,above,align=center},
  box/.style={draw=azulUNI!60,fill=white,rounded corners=1pt,font=\scriptsize,align=center,inner sep=3pt},
  reject/.style={draw=red!60,fill=red!6,rounded corners=1pt,font=\scriptsize,align=center,inner sep=3pt},
}

% Etiqueta para los pares de código
\newcommand{\codetag}[1]{\vspace{2pt}\noindent{\small\color{azulUNI}\textbf{#1}}\par\nobreak\vspace{1pt}}

\begin{document}

% ====================== PORTADA ======================
\begin{titlepage}
\centering
\vspace*{1.6cm}
{\color{azulUNI}\rule{\linewidth}{1.2pt}}\\[0.6cm]
{\LARGE\bfseries Del incidente a la recuperación\par}
\vspace{0.3cm}
{\Large Custodia de claves, cifrado híbrido y autenticación\par}
\vspace{0.2cm}
{\large\itshape Modelo A: el servidor genera y entrega la semilla maestra\par}
\vspace{0.4cm}
{\color{azulUNI}\rule{\linewidth}{1.2pt}}\\[1.2cm]

{\large\textbf{Curso:} Criptografía \quad\textbar\quad \textbf{Actividad III}\par}
\vspace{1.4cm}
{\large\textbf{Autores}\par}
\vspace{0.3cm}
{\large César Roa \\[0.2cm] Sebastián Ibáñez\par}

\vfill
{\large \today\par}
\end{titlepage}

% ====================== INTRODUCCIÓN ======================
\section{Introducción}

Para esta actividad elegimos el \textbf{Modelo A}: es el servidor quien genera la
semilla maestra $s_V$ de cada víctima y se la entrega ya cifrada. Nos pareció el
más natural para pensar la idea de \emph{custodia de claves}, porque deja muy
claro quién guarda el secreto (el servidor, como custodio) y quién lo usa sin
conservarlo (la víctima, que cifra sus archivos y luego olvida la semilla).

La idea completa es la siguiente. La víctima se registra ante el servidor y
recibe $s_V$ protegida con RSA. Con esa semilla deriva una clave distinta para
cada archivo y los cifra con AES. Después cierra todo y olvida la semilla; para
volver a abrir sus archivos debe \emph{recuperarla}, y hay dos caminos: pedirla al
servidor mediante un evento de pago, o reconstruirla por su cuenta investigando
el generador. A partir de ahí añadimos un atacante que se mete en medio de la
comunicación (MITM) y, finalmente, la defensa que lo neutraliza usando firmas
digitales.

A lo largo del informe explicamos qué tomamos del cuaderno de referencia (el
Colab 1) y, sobre todo, \textbf{qué tuvimos que cambiar y por qué}, mostrando los
fragmentos de código originales junto a los nuestros.

% ====================== CÓMO FUNCIONA ======================
\section{Cómo funciona nuestra aplicación}

El corazón del sistema es un módulo común, \texttt{crypto\_utils.py}, que reúne
todas las operaciones criptográficas, y dos programas que hablan por
\emph{sockets}: \texttt{servidor.py} y \texttt{cliente.py}. El flujo sigue los seis
pasos del enunciado.

Primero, cliente y servidor preparan su par de claves RSA-2048 y la víctima se
registra enviando su identificador $id_V$ y su clave pública $pk_V$; el servidor
responde con la suya. Después, el servidor genera la semilla con el generador
suministrado, la guarda asociada a esa víctima y se la manda cifrada con
RSA-OAEP. La víctima la descifra con su clave privada y la mantiene únicamente en
memoria. Con $s_V$ ya disponible, para cada archivo $i=0,1,2,\dots$ deriva una
clave
\[
K_i = \mathrm{HMAC\text{-}SHA256}\bigl(s_V,\ i\bigr),
\]
cifra el archivo con AES-256-GCM (nonce aleatorio de 12 bytes y \emph{tag} de 16)
y guarda todo en un pequeño archivo \texttt{.enc.json}. Un detalle importante es
que ligamos cada cifrado a su contexto mediante los datos asociados
$\mathrm{AAD} = id_V \parallel i$: así un archivo cifrado no puede ``colarse'' con
otro índice o de otra víctima sin que la verificación falle. Terminado el cifrado,
la víctima borra de memoria la semilla y las claves derivadas.

Para la recuperación por pago, la víctima vuelve a arrancar, carga de disco su
identidad y sus archivos cifrados, y pide la semilla con un evento de
autorización; el servidor, usando la clave pública que guardó en el registro,
le devuelve $s_V$ otra vez. Con ella regenera cada $K_i$, verifica el \emph{tag}
de GCM, guarda los archivos recuperados y los compara byte a byte contra los
originales. Decidimos que la entrega inicial ocurra \textbf{una sola vez}: si la
víctima intentara pedir de nuevo la semilla como ``primera entrega'', el servidor
la rechaza y la obliga a pasar por la vía de pago, que es justamente lo que el
enunciado busca proteger.

% ====================== COLAB ======================
\section{De dónde partimos: el Colab y cómo lo adaptamos}

Nos apoyamos en el \emph{Colab 1}, sobre todo en la sección~7 (RSA-OAEP) y la
sección~8 (la aplicación cliente-servidor). De ahí tomamos la estructura del
cifrado híbrido y la forma de comunicar los procesos. Ahora bien, el Colab usa la
biblioteca \texttt{PyCryptodome} y resuelve un problema algo distinto (establecer
una \emph{clave de sesión}), mientras que el enunciado nos pide usar
\texttt{cryptography} y trabajar con una \emph{semilla maestra} de la que se
derivan claves por archivo. Por eso, más que copiar, tuvimos que reescribir cada
pieza. A continuación mostramos los fragmentos originales y los nuestros.

\subsection{Transportar el secreto con RSA-OAEP}

En el Colab, el secreto aleatorio se encapsula y se recupera con OAEP así:

\codetag{Colab 1 (sección 8.2, PyCryptodome)}
\begin{lstlisting}
def encapsular_oaep(pk):
    secreto = secrets.token_bytes(32)
    publica = RSA.construct((pk["n"], pk["e"]))
    oaep = PKCS1_OAEP.new(publica, hashAlgo=SHA256)
    return oaep.encrypt(secreto), secreto

def desencapsular_oaep(sk, encapsulacion):
    privada = RSA.construct((sk["n"], sk["e"], sk["d"], sk["p"], sk["q"]))
    oaep = PKCS1_OAEP.new(privada, hashAlgo=SHA256)
    return oaep.decrypt(encapsulacion)
\end{lstlisting}

En nuestro proyecto hacemos lo mismo conceptualmente (cifrar y descifrar con
OAEP y SHA-256), pero con \texttt{cryptography}, que trabaja con objetos de clave
en lugar de tuplas de enteros:

\codetag{Nuestro código (\texttt{crypto\_utils.py}, cryptography)}
\begin{lstlisting}
_OAEP = padding.OAEP(mgf=padding.MGF1(hashes.SHA256()),
                     algorithm=hashes.SHA256(), label=None)

def rsa_cifrar(clave_publica, mensaje):
    return clave_publica.encrypt(mensaje, _OAEP)

def rsa_descifrar(clave_privada, cifrado):
    return clave_privada.decrypt(cifrado, _OAEP)
\end{lstlisting}

La operación criptográfica es la
misma (OAEP con MGF1 y SHA-256). Lo que cambia es la biblioteca y, con ella, el
manejo de las claves: el Colab reconstruye la clave desde sus componentes
(\texttt{RSA.construct}), mientras que nosotros las serializamos en formato PEM y
las cargamos como objetos. Además, en el Colab el secreto encapsulado es una
clave de sesión que se genera ahí mismo; en el nuestro, el ``secreto'' que viaja
dentro de OAEP es la semilla maestra $s_V$ que produce el generador suministrado.

\subsection{Derivar las claves: de HKDF a una clave por archivo}

Esta es la diferencia de fondo. El Colab deriva \emph{una} clave de sesión con
HKDF a partir del secreto y una sal:

\codetag{Colab 1 (sección 8.2, PyCryptodome)}
\begin{lstlisting}
def derivar_clave_sesion(secreto, sal):
    if len(sal) != 16:
        raise ValueError("La sal de sesion debe tener 16 bytes.")
    return HKDF(secreto, 32, sal, SHA256, context=b"Colab01/sesion")
\end{lstlisting}

Nosotros no queremos una única clave de sesión, sino una clave \emph{distinta por
archivo}, y además el enunciado fija exactamente la fórmula. Por eso usamos HMAC
con el contador del archivo como mensaje:

\codetag{Nuestro código (\texttt{crypto\_utils.py})}
\begin{lstlisting}
def derivar_clave(s_v, i):
    h = hmac.HMAC(s_v, hashes.SHA256())
    h.update(i.to_bytes(8, "big"))   # K_i = HMAC-SHA256(s_v, i)
    return h.finalize()
\end{lstlisting}

Nuestra razón para cambiarlo es sencilla. Ambos esquemas derivan claves a partir
de un secreto usando HMAC por debajo, pero con propósitos distintos. HKDF produce
una sola clave de sesión para toda la conexión; nosotros necesitamos que cada
archivo $i$ tenga su propia clave $K_i$, reproducible más tarde solo con la
semilla y el índice. Así, en la recuperación basta volver a calcular
$K_i=\mathrm{HMAC}(s_V,i)$ para cada archivo, sin guardar ninguna clave en disco.
Esta decisión es también la que hace posible el ataque de investigación (5b): como
$K_0$ depende solo de la semilla y de $i=0$, quien adivine la semilla reconstruye
la clave.

\subsection{Cifrar con AES-GCM}

El Colab cifra cada envío con AES-GCM y, para no repetir nonces, usa un contador
como nonce:

\codetag{Colab 1 (sección 8.3, PyCryptodome)}
\begin{lstlisting}
def cifrar_simetrico(clave, datos, contador):
    nonce = contador.to_bytes(12, "big")
    aes = AES.new(clave, AES.MODE_GCM, nonce=nonce, mac_len=16)
    cifrado, etiqueta = aes.encrypt_and_digest(datos)
    return nonce + etiqueta + cifrado
\end{lstlisting}

Nosotros también usamos AES-GCM con \emph{tag} de 16 bytes, pero tomamos el nonce
al azar y añadimos datos asociados que ligan el cifrado a la víctima y al índice:

\codetag{Nuestro código (\texttt{crypto\_utils.py})}
\begin{lstlisting}
def construir_aad(id_v, i):
    return id_v.encode("utf-8") + b"|" + i.to_bytes(8, "big")

def aes_cifrar(clave, plaintext, aad):
    nonce = os.urandom(12)
    ct = AESGCM(clave).encrypt(nonce, plaintext, aad)
    return nonce, ct[:-16], ct[-16:]      # nonce, ciphertext, tag
\end{lstlisting}

La diferencia aquí está en el nonce. El Colab lo deriva del contador
porque en una sesión corta eso basta para que no se repita y, de paso, le sirve
para detectar paquetes fuera de orden. En nuestro caso cada archivo se cifra con
una clave $K_i$ distinta y queremos guardarlos de forma independiente, así que un
nonce aleatorio es más simple y seguro (lo guardamos junto al cifrado). Donde el
Colab usa el contador como control anti-repetición, nosotros usamos los datos
asociados $id_V \parallel i$: no cifran nada, pero quedan ``firmados'' por el
\emph{tag}, de modo que si alguien cambia el identificador o el índice, el
descifrado falla.

\subsection{La comunicación cliente-servidor}

De la sección~8 también tomamos la forma de enviar mensajes por TCP. Como TCP es
un flujo de bytes sin fronteras, el Colab antepone a cada paquete su longitud en
cuatro bytes. Nosotros hacemos exactamente eso, pero enviando objetos JSON (lo que
nos da mensajes con campos claros como \texttt{accion} o \texttt{id\_v}):

\codetag{Nuestro código (\texttt{crypto\_utils.py})}
\begin{lstlisting}
def enviar_json(sock, obj):
    datos = json.dumps(obj).encode("utf-8")
    sock.sendall(struct.pack(">I", len(datos)) + datos)   # longitud(4B) || datos
\end{lstlisting}

En resumen, del Colab conservamos el cifrado híbrido (OAEP para el secreto y
AES-GCM para los datos) y el estilo de comunicación; cambiamos la biblioteca, la
naturaleza del secreto (una semilla maestra en vez de una clave de sesión), la
derivación (una clave por archivo en vez de una de sesión) y el manejo del nonce;
y añadimos toda la parte de persistencia y de autenticación, que el Colab da por
supuesta.

% ====================== FLUJOS Y DIAGRAMA ======================
\section{El protocolo con autenticación y el punto de rechazo}

El enunciado nos pide adaptar el diagrama del modelo añadiendo los datos firmados,
las verificaciones y el punto exacto donde se detiene un ataque. La
Figura~\ref{fig:legitimo} muestra el intercambio legítimo ya autenticado y la
Figura~\ref{fig:mitm}, el ataque del intermediario detenido por la defensa. Las
claves públicas de verificación están \textbf{pre-compartidas} en
\texttt{trusted\_keys/} y no viajan por la red.

\begin{figure}[h]\centering
\begin{tikzpicture}
  \node[actor] (V) at (0,0) {Víctima\\\footnotesize(Bob)};
  \node[actor] (S) at (9,0) {Servidor\\\footnotesize(Alice)};
  \draw[dashed,gray] (V) -- (0,-7.2);
  \draw[dashed,gray] (S) -- (9,-7.2);
  \draw[msg] (0,-1.0) -- node[lbl]{\textsf{REGISTRO:} $id_V,\ pk_V,$ firma$_C$} (9,-1.0);
  \draw[msg] (9,-1.8) -- node[lbl,below]{\textsf{REGISTRO\_OK:} $pk_S$} (0,-1.8);
  \node[box] at (9,-2.55) {Verifica firma$_C$ con\\ confianza del cliente};
  \draw[msg] (0,-3.2) -- node[lbl]{\textsf{SOLICITAR\_SEMILLA:} $id_V$} (9,-3.2);
  \node[box] at (9,-3.95) {$s_V \leftarrow$ generar(); entrega única};
  \draw[msg] (9,-4.6) -- node[lbl,below]{\textsf{SEMILLA:} $\mathrm{OAEP}(pk_V,s_V),$ firma$_S$} (0,-4.6);
  \node[box] at (0,-5.35) {Verifica firma$_S$\\ \emph{antes} de usar $s_V$};
  \draw[msg] (0,-6.0) -- node[lbl]{\textsf{ACK:} $id_V$} (9,-6.0);
  \node[box,align=center] at (4.5,-6.8)
    {Paso 3: $K_i=\mathrm{HMAC}(s_V,i)$; \ AES-GCM$(K_i,\ \mathrm{AAD}=id_V\parallel i)$; \ borra $s_V,K_i$};
\end{tikzpicture}
\caption{Intercambio legítimo autenticado. La recuperación por pago repite el
patrón: el cliente firma \texttt{RECUPERAR} y el servidor devuelve
$\mathrm{OAEP}(pk_V,s_V)$ firmado, que el cliente verifica antes de descifrar.}
\label{fig:legitimo}
\end{figure}

\begin{figure}[h]\centering
\begin{tikzpicture}
  \node[actor] (V) at (0,0) {Víctima};
  \node[actor] (M) at (4.5,0) {Proxy (MITM)};
  \node[actor] (S) at (9,0) {Servidor};
  \draw[dashed,gray] (V) -- (0,-5.4);
  \draw[dashed,gray] (M) -- (4.5,-5.4);
  \draw[dashed,gray] (S) -- (9,-5.4);
  \draw[msg] (0,-1.0) -- node[lbl]{$id_V, pk_V,$ firma$_C$} (4.5,-1.0);
  \draw[msg] (4.5,-1.9) -- node[lbl]{$id_V, \mathbf{pk_M},$ firma$_C$} (9,-1.9);
  \node[reject,align=center] at (9,-2.75)
    {\textbf{Caso 1:} Verify(firma$_C$ sobre $pk_M$) = FALLA};
  \node[reject,align=center] at (9,-3.55)
    {\textbf{Caso 2:} clave de firma adjunta $\notin$ \texttt{trusted\_keys}};
  \draw[rmsg,red!70] (9,-4.2) -- node[lbl,below]{\textsf{ABORTADO}} (4.5,-4.2);
  \draw[rmsg,red!70] (4.5,-4.2) -- node[lbl,below]{\textsf{ABORTADO}} (0,-4.2);
  \node[reject,align=center] at (4.5,-5.0)
    {El servidor nunca genera ni envía $s_V$: \textbf{rechazo en el paso 1}};
\end{tikzpicture}
\caption{El ataque MITM detenido por la defensa. El servidor aborta antes de
liberar la semilla porque valida exclusivamente contra \texttt{trusted\_keys/}.}
\label{fig:mitm}
\end{figure}

El rechazo ocurre en el \textbf{paso 1}, en el servidor, antes de generar o enviar
la semilla. Lo explicamos en detalle en la sección de la defensa.

% ====================== ENTROPÍA ======================
\section{Rompiendo la semilla: análisis de entropía}

La parte más interesante del trabajo es que, aun usando RSA-OAEP y AES-256-GCM
(primitivas fuertes), el sistema se puede romper por completo. El motivo no está
en el cifrado, sino en el generador suministrado:

\begin{lstlisting}
configuracion = secrets.randbelow(9) + 16    # c uniforme en {16,...,24}
valor         = secrets.randbelow(1 << configuracion)  # uniforme en [0, 2^c)
return valor.to_bytes(32, "big")
\end{lstlisting}

Aunque la semilla ocupa 32 bytes, su valor nunca supera $2^{24}$, porque
$c\le 24$. Es decir, el espacio real de semillas es $[0,2^{24})$: unos 16,7
millones de posibilidades, frente a los $2^{256}$ que uno esperaría. Ese espacio
se recorre por fuerza bruta en segundos.

\paragraph{La distribución no es uniforme.} Como $c$ se elige uniforme entre 16 y
24 y luego el valor es uniforme en $[0,2^c)$, los números pequeños aparecen bajo
más escalas y por tanto son más probables. Formalmente,
\[
\Pr[\texttt{valor}=v]=\frac{1}{9}\sum_{c=16}^{24}\frac{\mathbf{1}[\,v<2^{c}\,]}{2^{c}},
\qquad
\Pr[\texttt{valor}<2^{16}]=\frac{1}{9}\left(2-2^{-8}\right)\approx 0{,}22 .
\]
O sea, más de uno de cada cinco valores cae en el $0{,}4\%$ más bajo del rango.

\paragraph{Cuántos intentos esperamos.} Si buscamos desde $0$ hacia arriba, el
número de intentos es $N=\texttt{valor}+1$. Calculamos
\[
\mathbb{E}[2^{c}]=\frac{2^{25}-2^{16}}{9}\approx 3\,720\,988,
\qquad
\mathbb{E}[\texttt{valor}]\approx 1\,860\,494,
\qquad
\boxed{\ \mathbb{E}[N]\approx 1\,860\,495\ }.
\]
En el peor caso tendríamos que llegar a $2^{24}\approx 1{,}68\times10^{7}$. Para
encontrar la semilla solo necesitamos el archivo cifrado con $i=0$: por cada
candidato derivamos $K_0=\mathrm{HMAC}(\text{candidato},0)$ e intentamos descifrar;
el \emph{tag} de GCM nos dice cuándo acertamos.

\paragraph{Lo que medimos.} En una de nuestras corridas la semilla se encontró
tras $150\,072$ intentos en $1{,}42$ s, a unos $1{,}1\times10^5$ candidatos por
segundo; en otra bastaron $13\,619$ intentos en $0{,}10$ s. Las cifras varían
porque dependen de la semilla concreta, pero siempre quedan muy por debajo del
peor caso. A esa velocidad, incluso agotar todo el espacio tomaría alrededor de
dos minutos y medio.

\paragraph{Romper la semilla frente a descifrar los archivos.} Conviene separar
los dos costos. Recuperar $s_V$ cuesta del orden de $\mathbb{E}[N]\approx 1{,}9$
millones de operaciones; descifrar los $M$ archivos, una vez que ya tenemos la
semilla, cuesta solo $M$ operaciones. El trabajo duro es, con diferencia, romper
la semilla: para $M=3$ archivos la proporción es de unas $620\,000$ a~$1$. En la
práctica se paga una vez por quebrar la semilla y después descifrar cualquier
cantidad de archivos es casi gratis. Toda la fortaleza de RSA y AES queda anulada
por el eslabón débil, que es la entropía del generador.

% ====================== MITM ======================
\section{El ataque del intermediario (MITM)}

Para mostrar por qué hace falta autenticación, construimos un proxy que se sienta
entre la víctima y el servidor y reenvía los mensajes. En el registro, sustituye
la clave pública de la víctima $pk_V$ por la suya, $pk_M$. El servidor, que no
tiene forma de saber de quién es realmente esa clave, guarda $pk_M$ creyendo que
es de la víctima y cifra la semilla hacia ella. El proxy la descifra con su propia
clave privada, se queda con $s_V$ y luego la vuelve a cifrar con la $pk_V$
auténtica para reenviársela a la víctima. La víctima completa su cifrado sin notar
absolutamente nada, y el proxy ya tiene la semilla con la que descifrar todos los
archivos.

Comprobamos este ataque de principio a fin: el proxy registra la semilla
interceptada y descifra una copia de los archivos de la víctima, con el servidor
luego apagado. Funciona precisamente porque RSA-OAEP garantiza confidencialidad
hacia ``quien tenga la clave'', pero no dice nada sobre de quién es esa clave.

% ====================== DEFENSA ======================
\section{La defensa con firmas digitales}

La defensa consiste en autenticar el origen de cada clave y de cada mensaje
importante, usando firmas Ed25519 cuyas claves públicas de verificación están
pre-compartidas. Tomamos cuatro decisiones de diseño.

\paragraph{Qué esquema de firma.} Elegimos Ed25519 (sobre RSA-PSS) porque es
determinista (no depende de aleatoriedad al firmar, lo que evita un tipo de fallo
clásico), es rápido, produce firmas cortas de 64 bytes y su interfaz es difícil de
usar mal. Firmar y cifrar quedan como funciones separadas: RSA-OAEP sigue
protegiendo la semilla, y las firmas solo autentican.

\paragraph{Qué firmamos.} Firmamos siempre el trío
$\ id_V \parallel pk_{\text{cifrado}} \parallel \text{rol}$. El $id_V$ ata la firma
a una víctima concreta; $pk_{\text{cifrado}}$ es justo lo que viaja en claro en ese
paso (la clave pública en el registro, o el criptograma de la semilla en el paso
2), de modo que si el proxy lo altera la firma deja de cuadrar; y el rol
(``cliente'' o ``servidor'') evita que una firma se reutilice en la otra
dirección. Es la misma idea de los datos asociados del cifrado: atar cada dato a
su contexto.

\paragraph{De dónde sale la confianza.} Las claves públicas de verificación se
reparten por un canal aparte y se guardan en \texttt{trusted\_keys/}; no se envían
por la red. Son el ancla de confianza: las firmas no crean confianza, solo la
trasladan desde ese ancla. Por eso el servidor valida \emph{únicamente} contra esas
claves. Esto nos deja frenar los dos intentos del proxy: en el \textbf{Caso 1}
cambia la clave pero no la firma, y la verificación falla; en el \textbf{Caso 2}
forja una firma con su propia clave y adjunta su clave de verificación, pero como
esa clave no está en \texttt{trusted\_keys/}, el servidor la rechaza igual.

\paragraph{Cuándo se verifica.} La comprobación ocurre \emph{antes} de liberar la
semilla. El servidor verifica la firma del cliente al recibir el registro y, si no
cuadra, aborta sin generar ni enviar $s_V$. De forma simétrica, el cliente
verifica la firma del servidor sobre el criptograma antes de descifrar la semilla.
Comprobamos que, con el proxy en medio, el protocolo se detiene y la semilla nunca
se roba; y que, con los participantes legítimos, el ciclo completo (cifrado,
reinicio y recuperación por pago) funciona con normalidad.

% ====================== PRUEBAS ======================
\section{Qué probamos}

Para respaldar todo lo anterior dejamos un script, \texttt{ejecutar\_pruebas.sh},
que ejecuta las cuatro pruebas pedidas y guarda cada salida en \texttt{evidencias/}.
La primera prueba cifra los archivos, reinicia y recupera por pago: los tres
archivos vuelven idénticos a los originales (comparación byte a byte y por
SHA-256). La segunda es la investigación con el servidor apagado: encuentra la
semilla por fuerza bruta y descifra todo. La tercera es el MITM sin defensa: el
proxy intercepta la semilla y descifra el botín. La cuarta prueba la defensa:
rechaza el Caso 1, rechaza el Caso 2 y, con los participantes legítimos, completa
el ciclo con recuperación; en ninguno de los dos ataques el proxy logra quedarse
con la semilla.

También dejamos todo listo para correr en contenedores: un \texttt{Dockerfile}
ligero y un \texttt{docker-compose.yml} que aíslan servidor, proxy y cliente en
una red interna, cada uno con su propio volumen de datos.

% ====================== CONCLUSIÓN ======================
\section{Conclusión}

Nos quedamos con dos ideas. La primera es que un sistema puede usar criptografía
impecable y aun así ser inseguro si un solo componente, aquí el generador de la
semilla, tiene poca entropía; el cifrado no compensa ese fallo. La segunda es que
cifrar no es autenticar: RSA-OAEP protege la semilla, pero sin firmas ancladas en
una confianza previa un intermediario puede sustituir claves y quedarse con el
secreto. Ambas ideas se ven de forma concreta en el código y en las evidencias que
acompañan este informe.

\end{document}
